Throughout this thesis, bold symbols denote vectors and tensors, while non-bold symbols denote scalar quantities. When no component index is specified, the corresponding non-bold symbol denotes the magnitude of the bold quantity; for example,
\[
    r := |\mathbf{r}|.
\]
Relative quantities between nodes $i$ and $j$ are written with a double subscript, for example
\[
    \mathbf{r}_{ij} := \mathbf{r}_j - \mathbf{r}_i.
\]

We adopt the Einstein summation convention throughout, so repeated indices are implicitly summed over, and partial derivatives with respect to Cartesian coordinates are written as
\[
    \partial_i := \frac{\partial}{\partial x_i}.
\]
A circumflex (``hat'') denotes the Fourier transform of a quantity, so that for a function $f(x)$ we write
\[
    \hat{f}(k) := \mathcal{F}(f)(k),
\]
where $\mathcal{F}$ denotes the Fourier transform and $k$ is the wavenumber.
A dot notation denotes the time derivative, so that
\[
    \dot{f} := \frac{df}{dt}.
\]
